% Placeholder. See manuscript-planning.md for the plan.
A number of datasets for exploring future climate change under different boundary conditions is now available.
These datasets cover those included in CMIP7's ScenarioMIP exercise [TODO cite van Vuuren].
As an illustration of our ability to produce greenhouse gas forcings
for different futures, we have also described the forcings produced for the counterfactual VL-CF experiment.
For users, a full user guide is available at [TODO link].
As for most forcings, the latest version
and versions to use in CMIP7 experiments can be found at
% \newline \mbox{\url{https://input4mips-cvs.readthedocs.io/en/latest/}},
\url{https://input4mips-cvs.readthedocs.io/en/latest/},
with the specific page for greenhouse gas forcing being
\newline \mbox{\url{https://input4mips-cvs.readthedocs.io/en/latest/dataset-overviews/greenhouse-gas-concentrations/}}.
Further guidance on how to use forcings in each CMIP experiment can be found at
\newline \mbox{\url{https://wcrp-cmip.github.io/cmip7-guidance/docs/CMIP7/Experiment_set_up_and_Forcings/}}.
As noted in the historical dataset description [TODO cite],
the conventions for recording which forcings were used in a simulation are not yet settled,
so the guidance pages, rather than this paper, should be used as the primary source of information.

Compared to the CMIP6 experiments,
the CMIP7 experiments cover a narrower range of greenhouse gas concentrations,
most notably for CO_2.
The CMIP7 experiments also include consideration of a greater variety of potential futures.
Notably, there is inclusion of sustained, large-scale negative emissions in the LN experiment
which may lead to CO_2 concentrations dropping below pre-industrial values
(emissions-driven ESM simulations for this experiment will be particularly novel),
greater exploration of the difference between peak warming in the range of 1.5 to 2.0\textdegree~C via the VL, L and LN experiments
and a more extreme exploration of reversibility via the HL experiment.
All of the changes are driven by changes in the underlying experiment descriptions between CMIP6 and CMIP7.

The dataset described here has been available from ESGF since approximately mid-2026,
keeping as close to CMIP7's intended timeline as possible (Sect.~\ref{sec:output-dataset-requirements}).
This meant that many potential improvements remain unaddressed (Sect.~\ref{sec:discussion}).
We aim to provide updates in future that improve our and the general field of forcing production and provision.
The European Union Horizon project Quality-assured Updates of forcing Information for Climate Change Assessments (QUICCA)
is an example of one project which is supporting such updates.
% % Include if we get confirmation of C3S before submission
% The [C3S project], despite its focus on historical forcings,
% will also support the development of a suite of tools
% that will also support efforts to produce forcings for projection experiments.
